\emph{完备性定理}是\emph{可靠性定理}的逆命题。
它的一种形式是：若 $\Gamma \Entails !A$ 则 $\Gamma \Proves !A$；
另一种形式是：若 $\Gamma$ 是一致的，则它是可满足的。
我们证明了第二种形式（并从第二种形式推导出了第一种形式）。
该证明较为复杂，需要多个步骤。
我们从一致集合~$\Gamma$ 出发。
首先，我们添加无穷多个新的常项符号~$c_i$，以及形如 $\lexists[x][!A(x)] \lif !A(c)$ 的公式，其中扩充后的语言中每个带自由变元的公式~$!A(x)$ 都与一个新常项配对。
这样就得到了一个包含~$\Gamma$ 的\emph{饱和的}一致语句集。
它仍然是一致的。
接着，我们取这个集合并将其扩充为一个\emph{完全一致集}。
完全一致集具有良好的性质：对任何语句~$!A$，要么 $!A$ 在集合中，要么 $\lnot !A$ 在集合中（但绝不会二者皆在其中）。
由于我们是从饱和集出发的，现在我们就得到了一个包含~$\Gamma$ 的饱和、完全、一致的语句集~$\Gamma^*$。
现在，从这个集合可以定义一个结构~$\Struct{M}$，使得 $\Sat{M(\Gamma^*)}{!A}$ 当且仅当 $!A \in \Gamma^*$。
特别地，$\Sat{M(\Gamma^*)}{\Gamma}$，即 $\Gamma$ 是可满足的。
如果语言中包含 $=$，构造会稍复杂一些。

从完备性定理可以得出两个重要推论。
\emph{紧致性定理}指出：$\Gamma \Entails !A$ 当且仅当存在某个有限子集 $\Gamma_0 \subseteq \Gamma$ 使得 $\Gamma_0 \Entails !A$。
一个等价的表述是：$\Gamma$ 是可满足的当且仅当每个有限子集 $\Gamma_0 \subseteq \Gamma$ 都是可满足的。
紧致性定理可用于证明具有特定性质的结构的存在性。
例如，我们可以用它证明：对于每个具有任意大有穷模型的理论，都存在无穷模型。
这尤其意味着，\emph{有穷性}无法在一阶逻辑中表达。
第二个推论是\emph{Löwenheim-Skolem 定理}，它指出：每个可满足的 $\Gamma$ 都有一个可数模型。
这进而表明，\emph{不可数性}也无法在一阶逻辑中表达。